\documentclass[11pt]{article}

\usepackage[margin=1in]{geometry}
\usepackage{amsmath,amssymb,amsthm,mathtools}
\usepackage{booktabs,longtable,array}
\usepackage{enumitem}
\usepackage{hyperref}
\usepackage{microtype}
\usepackage{xcolor}

\hypersetup{colorlinks=true,linkcolor=blue,citecolor=blue,urlcolor=blue}
\allowdisplaybreaks

\newtheorem{proposition}{Proposition}
\newtheorem{lemma}{Lemma}
\newtheorem{definition}{Definition}
\newtheorem{assumption}{Assumption}
\theoremstyle{remark}
\newtheorem{remark}{Remark}

\newcommand{\R}{\mathbb R}
\newcommand{\dd}{\,\mathrm d}
\newcommand{\E}{\mathbb E}
\newcommand{\widehatell}{\widehat\ell}
\newcommand{\widehatZ}{\widehat Z}
\newcommand{\widehatq}{\widehat q}
\newcommand{\widehatp}{\widehat p}
\newcommand{\dotwidehatZ}{\dot{\widehat Z}}
\newcommand{\dotwidehatq}{\dot{\widehat q}}
\newcommand{\dotwidehatp}{\dot{\widehat p}}
\newcommand{\vecop}{\operatorname{vec}}
\newcommand{\cond}{\,|\,}

\title{A Chair-Readable Derivation of the Fixed-Branch Squared Tensor-Train Gradient}
\author{BayesFilter Zhao--Cui Gradient Note}
\date{2026-06-01}

\begin{document}
\maketitle

\begin{abstract}
This note explains the analytical gradient used by the fixed-branch squared
tensor-train filtering construction inspired by Zhao and Cui.  The purpose is
not to introduce a new filtering method.  The purpose is to make one derivative
claim understandable: after all approximation choices are frozen, the derivative
computed by the backward calculation is the derivative of the same approximate
normalizing-constant scalar computed by the forward calculation.  The exposition
starts with scalar normalizing constants, then squared approximations, then a
two-coordinate tensor-train warmup, then fixed linear solves, and only then the
full tensor-train filtering derivative.
\end{abstract}

\tableofcontents

\section{Motivation For The Panel Chair}

The filtering algorithm produces, at each time \(t\), an approximate
normalizing constant
\[
    \widehat Z_t(\beta)>0.
    \tag{1}
\]
The parameter \(\beta\) is the model parameter at which we evaluate the
approximate likelihood.  The total approximate log likelihood is
\[
    \widehatell_T(\beta)
    =
    \sum_{t=1}^{T}\log \widehat Z_t(\beta).
    \tag{2}
\]
For a chemist, the closest familiar object is a partition function.  A
partition function is an integral of a nonnegative weight.  Here
\(\widehat Z_t(\beta)\) is also an integral of a nonnegative weight:
\[
    \widehat Z_t(\beta)
    =
    \int \widehat q_t(z;\beta)\dd z.
    \tag{3}
\]
The variable \(z\) is a numerical coordinate, usually a transformed coordinate
on a fixed box.  The function \(\widehat q_t\) is not the exact Bayesian
density.  It is the squared tensor-train approximation used by the numerical
filter.

The panel cares about the gradient because parameter inference algorithms need
to know how the approximate log likelihood changes when \(\beta\) changes:
\[
    \nabla_\beta \widehatell_T(\beta).
    \tag{4}
\]
The delicate point is not the derivative of \(\log\).  The delicate point is
that \(\widehat Z_t\) is produced by a complicated approximation algorithm.  If
the approximation algorithm changes its grid, ranks, pivots, or scaling shifts
when \(\beta\) changes, then the two likelihood evaluations may not be
evaluations of the same mathematical formula.  A classical analytical
derivative can only differentiate one declared formula.

That is why this note uses a fixed branch.

\begin{center}
\fbox{\begin{minipage}{0.92\linewidth}
\textbf{Fixed branch.}  Before differentiating, freeze the computational
choices: coordinate domain, fitting points, basis functions, tensor-train
ranks, sweep order, ridge parameter, defensive density, stabilizing shifts, and
the rule for constructing every fitted core.  Do not freeze the fitted core
values themselves.  Those core values are outputs of the fixed fitting
equations, so they remain functions of \(\beta\).  Then differentiate the
scalar defined by the saved choices and by the core values obtained from the
same saved fitting equations.
\end{minipage}}
\end{center}

The goal is therefore:
\[
    \text{differentiate the scalar actually computed by the saved forward pass.}
    \tag{5}
\]
The goal is not:
\[
    \text{differentiate every possible adaptive version of the algorithm.}
    \tag{6}
\]

\section{Warmup 1: A Normalizing Constant}

Start with the simplest possible object.  Let
\[
    Z(\beta)=\int_\Omega q(z;\beta)\dd z,
    \qquad
    q(z;\beta)\ge 0,
    \qquad
    0<Z(\beta)<\infty.
    \tag{7}
\]
The normalized density is
\[
    p(z;\beta)=\frac{q(z;\beta)}{Z(\beta)}.
    \tag{8}
\]
The log normalizing constant is
\[
    L(\beta)=\log Z(\beta).
    \tag{9}
\]
Assume that \(q(z;\beta)\) is differentiable in \(\beta\), and that we may move
the derivative inside the integral:
\[
    \frac{\partial}{\partial\beta}
    \int_\Omega q(z;\beta)\dd z
    =
    \int_\Omega
    \frac{\partial q(z;\beta)}{\partial\beta}\dd z.
    \tag{10}
\]
A sufficient condition is that, near the parameter value of interest, the
absolute derivative is bounded by an integrable function:
\[
    \left|
    \frac{\partial q(z;\beta)}{\partial\beta}
    \right|
    \le m(z),
    \qquad
    \int_\Omega m(z)\dd z<\infty.
    \tag{11}
\]
Then the derivative of the log normalizer is
\[
\begin{aligned}
    \frac{\partial L(\beta)}{\partial\beta}
    &=
    \frac{\partial}{\partial\beta}\log Z(\beta)\\
    &=
    \frac{1}{Z(\beta)}
    \frac{\partial Z(\beta)}{\partial\beta}\\
    &=
    \frac{1}{Z(\beta)}
    \int_\Omega
    \frac{\partial q(z;\beta)}{\partial\beta}\dd z.
\end{aligned}
\tag{12}
\]
This is the first idea in the whole gradient.  Everything later is only a way
to compute the numerator and denominator in (12) for the tensor-train
approximation.

For the sequential filter, the same formula is used at every time:
\[
    \frac{\partial\widehatell_T(\beta)}{\partial\beta}
    =
    \sum_{t=1}^{T}
    \frac{1}{\widehat Z_t(\beta)}
    \frac{\partial\widehat Z_t(\beta)}{\partial\beta}.
    \tag{13}
\]

\section{Warmup 2: A Squared Approximation}

The fixed-branch squared tensor-train approximation has the form
\[
    \widehat q(z;\beta)
    =
    e^{-c}\phi(z;\beta)^2+\tau\lambda(z).
    \tag{14}
\]
Here:
\[
\begin{array}{ll}
c & \text{is a frozen stabilizing shift},\\
\tau>0 & \text{is a frozen defensive mass},\\
\lambda(z)\ge0 & \text{is a frozen reference density with }\int\lambda=1,\\
\phi(z;\beta) & \text{is the fitted square-root function.}
\end{array}
\tag{15}
\]
Because \(c,\tau,\lambda\) are frozen, the only object in (14) that changes
with \(\beta\) is \(\phi\).  Differentiating (14) gives
\[
\begin{aligned}
    \dot{\widehat q}(z;\beta)
    &=
    \frac{\partial}{\partial\beta}
    \left[
    e^{-c}\phi(z;\beta)^2+\tau\lambda(z)
    \right]\\
    &=
    e^{-c}
    \frac{\partial}{\partial\beta}\phi(z;\beta)^2
    +
    \frac{\partial}{\partial\beta}\tau\lambda(z)\\
    &=
    2e^{-c}\phi(z;\beta)\dot\phi(z;\beta).
\end{aligned}
\tag{16}
\]
The dot means \(\partial/\partial\beta\).

The corresponding normalizer is
\[
\begin{aligned}
    \widehat Z(\beta)
    &=
    \int \widehat q(z;\beta)\dd z\\
    &=
    e^{-c}\int \phi(z;\beta)^2\dd z
    +
    \tau\int\lambda(z)\dd z\\
    &=
    e^{-c}\int \phi(z;\beta)^2\dd z+\tau.
\end{aligned}
\tag{17}
\]
Using (16), its derivative is
\[
\begin{aligned}
    \dot{\widehat Z}(\beta)
    &=
    \int \dot{\widehat q}(z;\beta)\dd z\\
    &=
    2e^{-c}
    \int \phi(z;\beta)\dot\phi(z;\beta)\dd z.
\end{aligned}
\tag{18}
\]
Substitution into (12) gives
\[
    \frac{\partial}{\partial\beta}\log\widehat Z(\beta)
    =
    \frac{
    2e^{-c}\int \phi(z;\beta)\dot\phi(z;\beta)\dd z}
    {e^{-c}\int\phi(z;\beta)^2\dd z+\tau}.
    \tag{19}
\]

If \(c\), \(\tau\), or \(\lambda\) changed with \(\beta\), then (16) would have
extra terms:
\[
    \dot{\widehat q}
    =
    -\dot c\,e^{-c}\phi^2
    +
    2e^{-c}\phi\dot\phi
    +
    \dot\tau\,\lambda
    +
    \tau\dot\lambda.
    \tag{20}
\]
The fixed-branch derivative does not use (20).  It uses the fixed-branch
formula (16).  This is not a minor technicality.  It is the mathematical line
between differentiating one saved scalar and differentiating a changing
adaptive procedure.

\section{Warmup 3: Two Coordinates, Rank One}

Before a full tensor train, consider a two-coordinate rank-one function:
\[
    \phi(z_1,z_2;\beta)=h_1(z_1;\beta)h_2(z_2;\beta).
    \tag{21}
\]
The square integral is
\[
\begin{aligned}
    R(\beta)
    &=
    \int\!\!\int
    \phi(z_1,z_2;\beta)^2\dd z_1\dd z_2\\
    &=
    \int\!\!\int
    h_1(z_1;\beta)^2h_2(z_2;\beta)^2\dd z_1\dd z_2\\
    &=
    \left[\int h_1(z_1;\beta)^2\dd z_1\right]
    \left[\int h_2(z_2;\beta)^2\dd z_2\right].
\end{aligned}
\tag{22}
\]
Define
\[
    R_1(\beta)=\int h_1(z_1;\beta)^2\dd z_1,
    \qquad
    R_2(\beta)=\int h_2(z_2;\beta)^2\dd z_2.
    \tag{23}
\]
Then
\[
    R(\beta)=R_1(\beta)R_2(\beta).
    \tag{24}
\]
Differentiate by the product rule:
\[
    \dot R(\beta)=\dot R_1(\beta)R_2(\beta)+R_1(\beta)\dot R_2(\beta).
    \tag{25}
\]
Each one-dimensional derivative is
\[
    \dot R_1(\beta)=2\int h_1(z_1;\beta)\dot h_1(z_1;\beta)\dd z_1,
    \tag{26}
\]
\[
    \dot R_2(\beta)=2\int h_2(z_2;\beta)\dot h_2(z_2;\beta)\dd z_2.
    \tag{27}
\]
Combining (25)--(27),
\[
\begin{aligned}
    \dot R(\beta)
    &=
    2
    \left[\int h_1\dot h_1\dd z_1\right]
    \left[\int h_2^2\dd z_2\right]\\
    &\quad+
    2
    \left[\int h_1^2\dd z_1\right]
    \left[\int h_2\dot h_2\dd z_2\right].
\end{aligned}
\tag{28}
\]
This is the simplest version of the tensor-train mass derivative: when one
factor changes, keep the other factors fixed and sum the contributions.

\section{Warmup 4: Two Coordinates, Rank \texorpdfstring{\(R\)}{R}}

Now let
\[
    \phi(z_1,z_2;\beta)
    =
    \sum_{r=1}^{R}h_{1r}(z_1;\beta)h_{2r}(z_2;\beta).
    \tag{29}
\]
The square integral is
\[
\begin{aligned}
    R(\beta)
    &=
    \int\!\!\int
    \left[
    \sum_{r=1}^{R}h_{1r}(z_1)h_{2r}(z_2)
    \right]
    \left[
    \sum_{s=1}^{R}h_{1s}(z_1)h_{2s}(z_2)
    \right]
    \dd z_1\dd z_2\\
    &=
    \sum_{r=1}^{R}\sum_{s=1}^{R}
    \left[\int h_{1r}(z_1)h_{1s}(z_1)\dd z_1\right]
    \left[\int h_{2r}(z_2)h_{2s}(z_2)\dd z_2\right].
\end{aligned}
\tag{30}
\]
Define the two mass matrices
\[
    M_1[r,s]=\int h_{1r}(z_1)h_{1s}(z_1)\dd z_1,
    \tag{31}
\]
\[
    M_2[r,s]=\int h_{2r}(z_2)h_{2s}(z_2)\dd z_2.
    \tag{32}
\]
Then
\[
    R(\beta)=\sum_{r,s}M_1[r,s]M_2[r,s].
    \tag{33}
\]
This is the Frobenius inner product
\[
    R(\beta)=\langle M_1(\beta),M_2(\beta)\rangle_F.
    \tag{34}
\]
Differentiate:
\[
    \dot R
    =
    \langle \dot M_1,M_2\rangle_F
    +
    \langle M_1,\dot M_2\rangle_F.
    \tag{35}
\]
The derivative of a mass matrix entry is
\[
    \dot M_1[r,s]
    =
    \int
    \dot h_{1r}(z_1)h_{1s}(z_1)
    +
    h_{1r}(z_1)\dot h_{1s}(z_1)
    \dd z_1,
    \tag{36}
\]
and similarly for \(M_2\).  The full tensor-train mass recursion is only this
calculation repeated along a chain of coordinates.

\section{Warmup 5: A Fixed Linear Solve}

The fitted tensor-train cores are obtained from finite-dimensional least-squares
updates.  One such update can be written
\[
    N(\beta)g(\beta)=d(\beta).
    \tag{37}
\]
Here \(g\) is the vector of unknown core coefficients in the current update,
\(N\) is the regularized normal matrix, and \(d\) is the right-hand side.
Differentiate both sides:
\[
    \frac{\partial}{\partial\beta}\{N(\beta)g(\beta)\}
    =
    \frac{\partial d(\beta)}{\partial\beta}.
    \tag{38}
\]
Use the product rule:
\[
    \dot N(\beta)g(\beta)+N(\beta)\dot g(\beta)=\dot d(\beta).
    \tag{39}
\]
Move the first term to the right:
\[
    N(\beta)\dot g(\beta)=\dot d(\beta)-\dot N(\beta)g(\beta).
    \tag{40}
\]
If \(N(\beta)\) is nonsingular, then
\[
    \dot g(\beta)=
    N(\beta)^{-1}\left[\dot d(\beta)-\dot N(\beta)g(\beta)\right].
    \tag{41}
\]
In computation, one should not form the inverse.  One solves the linear system
(40) using the same numerical linear algebra used for the forward solve.

The ridge term is important because the normal matrix has the form
\[
    N=A^\top W A+\rho I,
    \qquad \rho>0.
    \tag{42}
\]
If \(W\) is positive semidefinite, then
\[
    v^\top N v
    =
    v^\top A^\top W A v+\rho v^\top v
    \ge
    \rho\|v\|^2
    \tag{43}
\]
for all \(v\).  Hence \(N\) is positive definite and nonsingular.  Without
nonsingularity, the derivative of the solve may not be well defined.

\section{The Fixed-Branch Forward Pass}

At time \(t\), the fixed branch constructs a nonnegative approximation
\[
    \widehat q_t(z;\beta)
    =
    e^{-c_t}\phi_t(z;\beta)^2+\tau_t\lambda_t(z)
    \tag{44}
\]
and its normalizer
\[
    \widehat Z_t(\beta)=\int \widehat q_t(z;\beta)\dd z.
    \tag{45}
\]
The branch is the collection of structural choices that define the scalar as a
function of \(\beta\).  These choices include the domain, fitting points, basis,
ranks, sweep order, shift, defensive term, and ridge parameter.  The branch does
not mean that the fitted numerical core values are held constant when
\(\beta\) changes.  Instead,
\[
    C_{t,k}(\beta)
    =
    \text{the core obtained by the saved fitting rule at parameter }\beta.
    \tag{45a}
\]
The derivative pass computes \(\dot C_{t,k}(\beta)\).  This distinction matters:
freezing the core values would usually make the normalizer nearly constant and
would not test the derivative of the fitted approximation.

\begin{center}
\begin{tabular}{p{0.27\linewidth}p{0.31\linewidth}p{0.31\linewidth}}
\toprule
Forward object & What is stored & Why it is needed later\\
\midrule
Domain map \(\Psi_t\) & Coordinate intervals and Jacobian & Evaluates the same transformed target at every derivative point\\
Fitting points \(z^{(j)}\) & Deterministic point array and weights & Defines the least-squares equations\\
Basis \(b_k\) & Basis family, degree, mass matrices & Evaluates cores and mass contractions\\
Ranks \(R_k\) & Fixed rank vector & Fixes core dimensions\\
Sweep order & Core update order and number of sweeps & Defines a finite composition of solve maps\\
Shift \(c_t\) & One frozen scalar & Defines the same scaled square-root target\\
Defensive term & \(\tau_t,\lambda_t\) & Ensures nonnegative density floor\\
Cores \(C_{t,k}(\beta)\) & Parameter-dependent fitted coefficient arrays & Evaluate \(\phi_t\) and \(\widehat q_t\)\\
Mass contractions & \(M_{t,>k}(\beta)\), \(M_{t,<k}(\beta)\) & Compute integrals and marginals\\
Normalizer & \(\widehat Z_t\) & Adds \(\log\widehat Z_t\) to the scalar\\
Carried filter & Numerator \(a_t\) and normalized \(\widehat p_t\) & Feeds the next time step\\
\bottomrule
\end{tabular}
\end{center}

The transformed target fitted at time \(t\) is
\[
    \widetilde q_t(z;\beta)
    =
    \widehat p_{t-1}(x_{t-1}(z);\beta)
    f_\beta(x_t(z)\mid x_{t-1}(z))
    g_\beta(y_t\mid x_t(z))
    J_t(z).
    \tag{46}
\]
The fitted square-root values are
\[
    y_j(\beta)
    =
    e^{c_t/2}\sqrt{\widetilde q_t(z^{(j)};\beta)}.
    \tag{47}
\]
The tensor train is
\[
    \phi_t(z;\beta)
    =
    H_{t,1}(z_1;\beta)H_{t,2}(z_2;\beta)\cdots H_{t,D}(z_D;\beta),
    \tag{48}
\]
where \(H_{t,k}(z_k)\in\R^{R_{k-1}\times R_k}\) and \(R_0=R_D=1\).

The retained numerator for the next filter is
\[
    a_t(z_t;\beta)
    =
    \int \widehat q_t(z_t,z_{t-1};\beta)\dd z_{t-1}.
    \tag{49}
\]
The retained normalized filter is
\[
    \widehat p_t(z_t;\beta)
    =
    \frac{a_t(z_t;\beta)}{\widehat Z_t(\beta)}.
    \tag{50}
\]
If the filter is reported in physical coordinates, multiply by the appropriate
change-of-variables factor.  The derivative logic is unchanged if the domain
map is frozen.

\section{The Fixed-Branch Derivative Pass}

The derivative pass computes the derivative of each forward object that depends
on \(\beta\).  It does not create a new branch.

\begin{center}
\begin{tabular}{p{0.32\linewidth}p{0.32\linewidth}p{0.26\linewidth}}
\toprule
Derivative object & Forward object differentiated & Formula source\\
\midrule
\(\dot{\widetilde q}_{t,j}\) & target value at fitting point & product rule\\
\(\dot y_j\) & square-root fit value & chain rule\\
\(\dot A_k,\dot N_k,\dot d_k\) & least-squares system & environment derivatives\\
\(\dot C_{t,k}\) & fitted core & fixed linear solve\\
\(\dot M_{t,>k}\) & mass contraction & product rule in core indices\\
\(\dot{\widehat Z}_t\) & normalizer & squared-integral derivative\\
\(\dot a_t\) & carried numerator & marginal contraction derivative\\
\(\dot{\widehat p}_t\) & normalized carried filter & quotient rule\\
\(\dot{\widetilde q}_{t+1}\) & next target & product rule using \(\dot{\widehat p}_t\)\\
\bottomrule
\end{tabular}
\end{center}

\subsection{Target And Square-Root Target}

Differentiate (46).  At fitting point \(z^{(j)}\), abbreviate
\[
    \widehat p_{t-1,j}=\widehat p_{t-1}(x_{t-1}^{(j)};\beta),
    \quad
    f_j=f_\beta(x_t^{(j)}\mid x_{t-1}^{(j)}),
    \quad
    g_j=g_\beta(y_t\mid x_t^{(j)}).
    \tag{51}
\]
With \(J_{t,j}\) frozen, the derivative is
\[
\begin{aligned}
    \dot{\widetilde q}_{t,j}
    &=
    J_{t,j}
    \left[
    \dot{\widehat p}_{t-1,j} f_j g_j
    +
    \widehat p_{t-1,j}\dot f_j g_j
    +
    \widehat p_{t-1,j} f_j\dot g_j
    \right].
\end{aligned}
\tag{52}
\]
If the domain map depended on \(\beta\), an additional
\(\widehat p f g\,\dot J\) term would appear.  That is outside the fixed branch
used here.

From (47),
\[
    \dot y_j
    =
    e^{c_t/2}
    \frac{1}{2\sqrt{\widetilde q_{t,j}}}
    \dot{\widetilde q}_{t,j},
    \tag{53}
\]
provided \(\widetilde q_{t,j}>0\).  A numerical implementation should use a
positive floor if the target is close to zero.  If such a floor is used, the
flooring rule becomes part of the declared scalar.  For example,
\[
    y_j(\beta)
    =
    e^{c_t/2}\sqrt{\max\{\widetilde q_{t,j}(\beta),\varepsilon_{\rm floor}\}}
    \tag{53a}
\]
has derivative zero through the floor whenever
\(\widetilde q_{t,j}<\varepsilon_{\rm floor}\).  The formulas below describe
the smooth positive-target case; a floored implementation must differentiate
the floored scalar it actually evaluates.

\subsection{Design Matrix And Core Solve}

The design row is the first place where tensor-train notation can look more
mysterious than it is.  Fix one sample point \(z^{(j)}\) and one core \(k\).
All cores except core \(k\) are treated as the known left and right
environments for this local least-squares update.  Write
\[
    L_{j,k}[a]
    =
    \left[
    H_{1}(z_1^{(j)})\cdots H_{k-1}(z_{k-1}^{(j)})
    \right]_{1a},
    \tag{53b}
\]
\[
    R_{j,k}[b]
    =
    \left[
    H_{k+1}(z_{k+1}^{(j)})\cdots H_D(z_D^{(j)})
    \right]_{b1}.
    \tag{53c}
\]
The local core evaluated at \(z_k^{(j)}\) is
\[
    H_k(z_k^{(j)})[a,b]
    =
    \sum_{\ell=1}^{m_k}
    C_k[a,\ell,b]\,b_k(z_k^{(j)})[\ell],
    \tag{53d}
\]
where \(m_k\) is the number of one-dimensional basis functions in coordinate
\(k\).  The value of the tensor train at this one point is therefore
\[
\begin{aligned}
    \phi_t(z^{(j)};\beta)
    &=
    \sum_{a=1}^{R_{k-1}}
    \sum_{b=1}^{R_k}
    L_{j,k}[a]\,H_k(z_k^{(j)})[a,b]\,R_{j,k}[b]\\
    &=
    \sum_{a=1}^{R_{k-1}}
    \sum_{\ell=1}^{m_k}
    \sum_{b=1}^{R_k}
    L_{j,k}[a]\,
    b_k(z_k^{(j)})[\ell]\,
    R_{j,k}[b]\,
    C_k[a,\ell,b].
\end{aligned}
    \tag{53e}
\]
Equation (53e) is the whole reason the local fit is a linear least-squares
problem.  At the fixed sample point, the unknown numbers are the entries of
the core \(C_k\).  Every coefficient multiplying \(C_k[a,\ell,b]\) is already
known from the branch and from the other cores.

If \(g_k=\vecop(C_k)\) stacks entries in the order \((a,\ell,b)\), the row
coefficient of \(g_k[a,\ell,b]\) is
\[
    A_{j,k}[a,\ell,b]
    =
    L_{j,k}[a]\,
    b_k(z_k^{(j)})[\ell]\,
    R_{j,k}[b].
    \tag{53f}
\]
The compact Kronecker expression in (57) is just (53f) written without three
sums:
\[
    A_{j,k}
    =
    R_{j,k}^{\top}\otimes b_k(z_k^{(j)})^\top\otimes L_{j,k}.
    \tag{53g}
\]
To differentiate the row, keep the basis value fixed and differentiate the
two environments:
\[
\begin{aligned}
    \dot A_{j,k}[a,\ell,b]
    &=
    \dot L_{j,k}[a]\,
    b_k(z_k^{(j)})[\ell]\,
    R_{j,k}[b]\\
    &\quad+
    L_{j,k}[a]\,
    b_k(z_k^{(j)})[\ell]\,
    \dot R_{j,k}[b].
\end{aligned}
    \tag{53h}
\]
Equation (58) is (53h) written in the same Kronecker shorthand.  If the basis
or fitting point moved with \(\beta\), there would be an additional term with
\(\dot b_k(z_k^{(j)})\).  That term is absent because the branch keeps basis
and fitting points fixed.

For one core \(k\), the least-squares relation has rows
\[
    \phi_t(z^{(j)};\beta)=A_{j,k}(\beta)g_k(\beta),
    \tag{54}
\]
where \(g_k=\vecop(C_{t,k})\).  Let
\[
    L_{j,k}
    =
    H_{1}(z_1^{(j)})\cdots H_{k-1}(z_{k-1}^{(j)}),
    \tag{55}
\]
\[
    R_{j,k}
    =
    H_{k+1}(z_{k+1}^{(j)})\cdots H_D(z_D^{(j)}).
    \tag{56}
\]
Then
\[
    A_{j,k}
    =
    R_{j,k}^{\top}\otimes b_k(z_k^{(j)})^\top\otimes L_{j,k}.
    \tag{57}
\]
Because basis values and fitting points are frozen,
\[
    \dot A_{j,k}
    =
    \dot R_{j,k}^{\top}\otimes b_k(z_k^{(j)})^\top\otimes L_{j,k}
    +
    R_{j,k}^{\top}\otimes b_k(z_k^{(j)})^\top\otimes \dot L_{j,k}.
    \tag{58}
\]
The environments are differentiated by product rule:
\[
    \dot L_{j,1}=0,
    \qquad
    \dot L_{j,k+1}
    =
    \dot L_{j,k}H_k(z_k^{(j)})
    +
    L_{j,k}\dot H_k(z_k^{(j)}),
    \tag{59}
\]
\[
    \dot R_{j,D}=0,
    \qquad
    \dot R_{j,k-1}
    =
    \dot H_k(z_k^{(j)})R_{j,k}
    +
    H_k(z_k^{(j)})\dot R_{j,k}.
    \tag{60}
\]
The local core matrix derivative is
\[
    \dot H_k(z_k^{(j)})[a,b]
    =
    \sum_{\ell}
    \dot C_k[a,\ell,b]\,b_k(z_k^{(j)})[\ell].
    \tag{61}
\]

The fixed ridge solve is
\[
    N_k g_k=d_k,
    \qquad
    N_k=A_k^\top W A_k+\rho I,
    \qquad
    d_k=A_k^\top W y.
    \tag{62}
\]
Using (40),
\[
    N_k\dot g_k=\dot d_k-\dot N_k g_k,
    \tag{63}
\]
where
\[
    \dot N_k=\dot A_k^\top W A_k+A_k^\top W\dot A_k,
    \tag{64}
\]
\[
    \dot d_k=\dot A_k^\top W y+A_k^\top W\dot y.
    \tag{65}
\]
The derivative update is therefore another linear solve with the same
left-hand matrix \(N_k\).  This is why the fixed branch is implementable: the
forward pass and derivative pass use the same stored computational path.

\subsection{Mass Contraction And Normalizer}

The full right-mass recursion is the same idea as the rank-\(R\) warmup in
(29)--(36).  In the warmup, \(M_1[r,s]\) stored the inner products of the
first-coordinate factors and \(M_2[r,s]\) stored the inner products of the
second-coordinate factors.  In the full tensor train, \(M_{>j}[b,b']\) plays
the role of the already-integrated right factor, \(B_j[\ell,\ell']\) is the
one-coordinate basis inner product, and the two copies of \(C_j\) are the two
factors coming from squaring \(\phi_t\).

The right mass contraction is
\[
    M_{>j-1}[a,a']
    =
    \sum_{b,b',\ell,\ell'}
    C_j[a,\ell,b]
    M_{>j}[b,b']
    C_j[a',\ell',b']
    B_j[\ell,\ell'].
    \tag{66}
\]
Here \(B_j\) is the basis mass matrix.  It is frozen.  Differentiate (66):
\[
\begin{aligned}
    \dot M_{>j-1}[a,a']
    &=
    \sum_{b,b',\ell,\ell'}
    \dot C_j[a,\ell,b]
    M_{>j}[b,b']
    C_j[a',\ell',b']
    B_j[\ell,\ell']\\
    &\quad+
    \sum_{b,b',\ell,\ell'}
    C_j[a,\ell,b]
    \dot M_{>j}[b,b']
    C_j[a',\ell',b']
    B_j[\ell,\ell']\\
    &\quad+
    \sum_{b,b',\ell,\ell'}
    C_j[a,\ell,b]
    M_{>j}[b,b']
    \dot C_j[a',\ell',b']
    B_j[\ell,\ell'].
\end{aligned}
\tag{67}
\]
Start from
\[
    M_{>D}=1,
    \qquad
    \dot M_{>D}=0.
    \tag{68}
\]
Then (66)--(67) move from the right end of the chain to the left end.  The
final square integral is
\[
    R_t(\beta)=\int\phi_t(z;\beta)^2\dd z=M_{>0}.
    \tag{69}
\]
Its derivative is
\[
    \dot R_t(\beta)=\dot M_{>0}.
    \tag{70}
\]
By (17)--(18),
\[
    \widehat Z_t(\beta)=e^{-c_t}R_t(\beta)+\tau_t,
    \tag{71}
\]
\[
    \dot{\widehat Z}_t(\beta)=e^{-c_t}\dot R_t(\beta).
    \tag{72}
\]

\subsection{Carried Filter And Next Time Step}

The carried numerator is
\[
    a_t(z_t;\beta)
    =
    \int \widehat q_t(z_t,z_{t-1};\beta)\dd z_{t-1}.
    \tag{73}
\]
Its derivative is
\[
    \dot a_t(z_t;\beta)
    =
    \int \dot{\widehat q}_t(z_t,z_{t-1};\beta)\dd z_{t-1}.
    \tag{74}
\]
For the squared part of the approximation, this marginal is another tensor
contraction.  To make the formula concrete, suppose the retained coordinate is
the last coordinate \(z_D=z_t\), and the previous-state coordinates are
\(z_1,\ldots,z_{D-1}\).  First integrate out the previous-state coordinates by
the same left-to-right mass recursion:
\[
    M_{<0}=1,
    \qquad
    M_{<j}[b,b']
    =
    \sum_{a,a',\ell,\ell'}
    M_{<j-1}[a,a']\,
    C_j[a,\ell,b]\,
    C_j[a',\ell',b']\,
    B_j[\ell,\ell']
    \tag{74a}
\]
for \(j=1,\ldots,D-1\).  The derivative of this marginal environment is
\[
\begin{aligned}
    \dot M_{<j}[b,b']
    &=
    \sum_{a,a',\ell,\ell'}
    \dot M_{<j-1}[a,a']\,
    C_j[a,\ell,b]\,
    C_j[a',\ell',b']\,
    B_j[\ell,\ell']\\
    &\quad+
    \sum_{a,a',\ell,\ell'}
    M_{<j-1}[a,a']\,
    \dot C_j[a,\ell,b]\,
    C_j[a',\ell',b']\,
    B_j[\ell,\ell']\\
    &\quad+
    \sum_{a,a',\ell,\ell'}
    M_{<j-1}[a,a']\,
    C_j[a,\ell,b]\,
    \dot C_j[a',\ell',b']\,
    B_j[\ell,\ell'].
\end{aligned}
    \tag{74b}
\]
Start with \(\dot M_{<0}=0\).  After the recursion reaches \(D-1\), the
retained-coordinate numerator is the function
\[
\begin{aligned}
    a_t(z_D;\beta)
    &=
    e^{-c_t}
    \sum_{b,b',\ell,\ell'}
    M_{<D-1}[b,b']\,
    C_D[b,\ell,1]\,
    C_D[b',\ell',1]\,
    b_D(z_D)[\ell]\,
    b_D(z_D)[\ell']\\
    &\quad+
    \tau_t
    \int \lambda_t(z_D,z_{1:D-1})\dd z_{1:D-1}.
\end{aligned}
    \tag{74c}
\]
Its derivative, with \(c_t,\tau_t,\lambda_t\) and the basis fixed, is
\[
\begin{aligned}
    \dot a_t(z_D;\beta)
    &=
    e^{-c_t}
    \sum_{b,b',\ell,\ell'}
    \dot M_{<D-1}[b,b']\,
    C_D[b,\ell,1]\,
    C_D[b',\ell',1]\,
    b_D(z_D)[\ell]\,
    b_D(z_D)[\ell']\\
    &\quad+
    e^{-c_t}
    \sum_{b,b',\ell,\ell'}
    M_{<D-1}[b,b']\,
    \dot C_D[b,\ell,1]\,
    C_D[b',\ell',1]\,
    b_D(z_D)[\ell]\,
    b_D(z_D)[\ell']\\
    &\quad+
    e^{-c_t}
    \sum_{b,b',\ell,\ell'}
    M_{<D-1}[b,b']\,
    C_D[b,\ell,1]\,
    \dot C_D[b',\ell',1]\,
    b_D(z_D)[\ell]\,
    b_D(z_D)[\ell'].
\end{aligned}
    \tag{74d}
\]
Equations (74a)--(74d) are the coding recipe behind the short integral
identity (74).  If the retained state is not the last coordinate, reorder the
coordinates or use left and right environments around the retained block; the
same product-rule pattern applies.

The normalized carried filter is
\[
    \widehat p_t(z_t;\beta)
    =
    \frac{a_t(z_t;\beta)}{\widehat Z_t(\beta)}.
    \tag{75}
\]
Differentiate by quotient rule:
\[
    \dot{\widehat p}_t(z_t;\beta)
    =
    \frac{\dot a_t(z_t;\beta)}{\widehat Z_t(\beta)}
    -
    \frac{a_t(z_t;\beta)\dot{\widehat Z}_t(\beta)}
    {\widehat Z_t(\beta)^2}.
    \tag{76}
\]
This is the crucial bridge between time steps.  At time \(t+1\), the target
contains \(\widehat p_t\), so its derivative contains
\(\dot{\widehat p}_t\).  Without saving the carried-filter derivative, the
next-step gradient is incomplete.

\section{Proposition 1: The Fixed-Branch Recursion Is A Normalized Approximate Filter}

\begin{proposition}
Fix a parameter value \(\beta\), a time horizon \(T\), and fixed branch objects
for all times.  Suppose that each approximate joint density
\[
    \widehat q_t(z_t,z_{t-1};\beta)
    =
    e^{-c_t}\phi_t(z_t,z_{t-1};\beta)^2
    +
    \tau_t\lambda_t(z_t,z_{t-1})
    \tag{77}
\]
is nonnegative and integrable, and that
\[
    0<\widehat Z_t(\beta)
    =
    \int \widehat q_t(z_t,z_{t-1};\beta)\dd z_t\dd z_{t-1}
    <\infty.
    \tag{78}
\]
Define
\[
    \widehat p_t(z_t;\beta)
    =
    \frac{\int \widehat q_t(z_t,z_{t-1};\beta)\dd z_{t-1}}
    {\widehat Z_t(\beta)}.
    \tag{79}
\]
Then \(\widehat p_t(\cdot;\beta)\) is a normalized density in \(z_t\).  The
recursion is a Bayesian filtering recursion for the declared approximate joint
density \(\widehat q_t\), not a proof that \(\widehat q_t\) equals the exact
Bayesian joint density.
\end{proposition}

\begin{proof}
Nonnegativity follows immediately:
\[
    e^{-c_t}\phi_t^2\ge0,
    \qquad
    \tau_t\lambda_t\ge0.
    \tag{80}
\]
Therefore \(\widehat q_t\ge0\).  By assumption,
\(\widehat Z_t\) is positive and finite, so
\[
    \widehat p_t^{\rm joint}(z_t,z_{t-1};\beta)
    =
    \frac{\widehat q_t(z_t,z_{t-1};\beta)}
    {\widehat Z_t(\beta)}
    \tag{81}
\]
is a normalized joint density:
\[
\begin{aligned}
    \int\!\!\int
    \widehat p_t^{\rm joint}(z_t,z_{t-1};\beta)
    \dd z_t\dd z_{t-1}
    &=
    \frac{1}{\widehat Z_t(\beta)}
    \int\!\!\int
    \widehat q_t(z_t,z_{t-1};\beta)
    \dd z_t\dd z_{t-1}\\
    &=
    1.
\end{aligned}
\tag{82}
\]
The retained filter is the marginal of this normalized joint density:
\[
    \widehat p_t(z_t;\beta)
    =
    \int
    \widehat p_t^{\rm joint}(z_t,z_{t-1};\beta)
    \dd z_{t-1}.
    \tag{83}
\]
Integrating (83) over \(z_t\),
\[
\begin{aligned}
    \int \widehat p_t(z_t;\beta)\dd z_t
    &=
    \int
    \left[
    \int \widehat p_t^{\rm joint}(z_t,z_{t-1};\beta)
    \dd z_{t-1}
    \right]\dd z_t\\
    &=
    \int\!\!\int
    \widehat p_t^{\rm joint}(z_t,z_{t-1};\beta)
    \dd z_t\dd z_{t-1}\\
    &=
    1.
\end{aligned}
\tag{84}
\]
Thus \(\widehat p_t\) is normalized.  The proof used only the declared
approximation \(\widehat q_t\).  It did not use the exact Bayesian density
except as motivation for constructing \(\widehat q_t\).  Hence the theorem is a
normalization theorem for the approximate filtering recursion.
\end{proof}

\section{Proposition 2: The Fixed-Branch Derivative Differentiates The Same Scalar}

\begin{proposition}
Fix the branch \(B=(B_1,\ldots,B_T)\).  Assume:
\[
    \widehatell_T(\beta;B)=\sum_{t=1}^T\log\widehat Z_t(\beta;B_t),
    \tag{85}
\]
each \(\widehat Z_t\) is positive and finite, all model density evaluations are
differentiable in \(\beta\), every least-squares update uses a fixed
nonsingular ridge system, and differentiation under the finite-dimensional
contractions and integrals is valid.  Then the derivative pass defined by
(52)--(76) computes
\[
    \frac{\partial}{\partial\beta}
    \widehatell_T(\beta;B).
    \tag{86}
\]
\end{proposition}

\begin{proof}
The proof follows the computational graph of the fixed branch.

First, the transformed target (46) is a product of the previous carried filter,
transition density, observation density, and frozen Jacobian.  The product rule
therefore gives (52).  If \(t=1\), the previous filter is the initial density;
for \(t>1\), its derivative is supplied by the previous carried-filter
derivative (76).

Second, the fitted target values (47) are differentiable functions of
\(\widetilde q_{t,j}\).  The chain rule gives (53).  The stabilizing shift
\(c_t\) is frozen, so no \(\dot c_t\) term appears.

Third, each core update is a fixed linear solve.  Equations (37)--(41) prove
that differentiating the solve yields
\[
    N_k\dot g_k=\dot d_k-\dot N_k g_k.
    \tag{87}
\]
Equations (54)--(65) identify \(N_k,d_k,\dot N_k,\dot d_k\) for the fixed
least-squares update.  Since the number and order of solves are fixed, a finite
composition of differentiable solve maps gives differentiable fitted cores.

Fourth, the square integral \(R_t=\int\phi_t^2\) is computed by the mass
recursion (66).  Its derivative is the product-rule recursion (67).  Therefore
\[
    \dot R_t=\frac{\partial}{\partial\beta}
    \int\phi_t(z;\beta)^2\dd z.
    \tag{88}
\]
With frozen \(c_t,\tau_t,\lambda_t\), (71)--(72) give the derivative of
\(\widehat Z_t\).

Fifth, the carried numerator and carried filter are defined by (73)--(75).
Their derivatives are (74) and (76), obtained by differentiating under the
integral and by the quotient rule.  This supplies the derivative object needed
by the next time step.

Finally, since each time contribution is \(\log\widehat Z_t(\beta;B_t)\),
Warmup 1 gives
\[
    \frac{\partial}{\partial\beta}
    \log\widehat Z_t(\beta;B_t)
    =
    \frac{\dot{\widehat Z}_t(\beta;B_t)}
    {\widehat Z_t(\beta;B_t)}.
    \tag{89}
\]
Summing (89) over \(t=1,\ldots,T\) gives (86).

Every object differentiated in this proof is an object in the saved branch
\(B\).  The proof never differentiates a changed rank, changed grid, changed
shift, changed pivot set, or changed domain.  Therefore the derivative is the
derivative of the same scalar \(\widehatell_T(\beta;B)\), not the derivative of
an adaptive algorithm that chooses a new \(B\) when \(\beta\) changes.
\end{proof}

\section{What This Does Not Prove}

The result above is deliberately narrow.

It does not prove exact posterior accuracy:
\[
    \widehat q_t=q_t
    \quad\text{is not claimed.}
    \tag{90}
\]
It proves normalization and differentiation of the declared approximation.

It does not prove ranks stay small:
\[
    R_k\ \text{moderate for all }k
    \quad\text{is a diagnostic requirement, not a theorem here.}
    \tag{91}
\]
It does not differentiate adaptive decisions:
\[
    \frac{\partial}{\partial\beta}
    \{\text{chosen ranks, pivots, domains, shifts, fitting points}\}
    \quad\text{is not part of the formula.}
    \tag{92}
\]
It does not prove HMC convergence:
\[
    \text{correct local gradient of an approximate scalar}
    \ne
    \text{posterior sampling guarantee.}
    \tag{93}
\]

These caveats make the result more honest, not weaker.  The fixed-branch
gradient is a precise mathematical contract.  Empirical accuracy, rank
stability, and sampling performance must be tested separately.

\section{Finite-Difference Test}

The finite-difference test checks that the value path and derivative path agree
for the same branch.  Choose a base parameter \(\beta_0\).  Build the branch
once and save the structural choices
\[
    \begin{aligned}
    B=\{&
    \text{domains, points, basis, ranks, sweeps, shifts,}\\
    &\text{defensive terms, ridge parameters,}\\
    &\text{deterministic initialization rules}\}.
    \end{aligned}
    \tag{94}
\]
Then evaluate the same scalar at perturbed parameters while reusing those
structural choices and recomputing the parameter-dependent fitted values by the
same fixed equations:
\[
    \widehatell_T(\beta_0+h e_i;B),
    \qquad
    \widehatell_T(\beta_0-h e_i;B).
    \tag{95}
\]
That means the domains, points, ranks, shifts, and sweep order are unchanged,
but the fitted core values
\[
    C_{t,k}(\beta_0+h e_i;B)
    \quad\text{and}\quad
    C_{t,k}(\beta_0-h e_i;B)
    \tag{95a}
\]
are recomputed from the same saved fitting rule.  They are not copied from
\(\beta_0\).
The centered finite difference is
\[
    D_i(h)
    =
    \frac{
    \widehatell_T(\beta_0+h e_i;B)
    -
    \widehatell_T(\beta_0-h e_i;B)}
    {2h}.
    \tag{96}
\]
Compare it to the analytical derivative
\[
    G_i=\partial_i\widehatell_T(\beta_0;B).
    \tag{97}
\]
Record
\[
    e_i(h)=|D_i(h)-G_i|,
    \tag{98}
\]
and
\[
    e_i^{\rm rel}(h)
    =
    \frac{|D_i(h)-G_i|}{1+|G_i|}.
    \tag{99}
\]
Use a small schedule such as
\[
    h\in\{10^{-2},10^{-3},10^{-4},10^{-5}\}.
    \tag{100}
\]
A healthy derivative check typically shows decreasing error over a range of
\(h\), followed by roundoff or fitting-noise limitations.  If the error does
not decrease on any reasonable range, the derivative path is not yet aligned
with the value path.

The branch identity check is mandatory:
\[
    B(\beta_0)=B(\beta_0+h e_i)=B(\beta_0-h e_i)=B.
    \tag{101}
\]
If the implementation rebuilds ranks, points, shifts, or domains at the
perturbed values, then (96) is no longer testing the derivative of
\(\widehatell_T(\beta;B)\).  It is testing a different adaptive calculation.
If the implementation copies the numerical cores from \(\beta_0\) without
resolving the fixed fitting equations at the perturbed parameter, then it is
also testing the wrong object.  The correct test freezes the branch choices but
allows all differentiable outputs of those choices to change with \(\beta\).

\section{Minimal Runnable Example}

The smallest useful example should be nonlinear but one-dimensional:
\[
    x_t=\beta x_{t-1}+\sigma\epsilon_t,
    \qquad
    \epsilon_t\sim N(0,1),
    \tag{102}
\]
\[
    y_t=\sin(x_t)+\eta_t,
    \qquad
    \eta_t\sim N(0,r^2).
    \tag{103}
\]
Use \(T=2\).  The first time step builds
\[
    q_1(x_1,x_0;\beta)
    =
    p_\beta(x_0)
    f_\beta(x_1\mid x_0)
    g_\beta(y_1\mid x_1).
    \tag{104}
\]
It fits \(\phi_1\), computes
\[
    \widehat Z_1(\beta)=e^{-c_1}\int \phi_1(z_1,z_0;\beta)^2
    \dd z_1\dd z_0+\tau_1,
    \tag{105}
\]
and stores
\[
    \widehat p_1(z_1;\beta)
    =
    \frac{\int
    \left[e^{-c_1}\phi_1(z_1,z_0;\beta)^2+\tau_1\lambda_1(z_1,z_0)\right]
    \dd z_0}
    {\widehat Z_1(\beta)}.
    \tag{106}
\]
The second time step must call that stored object:
\[
    q_2(x_2,x_1;\beta)
    =
    \widehat p_1(x_1;\beta)
    f_\beta(x_2\mid x_1)
    g_\beta(y_2\mid x_2).
    \tag{107}
\]
The scalar printed by the example is
\[
    \widehatell_2(\beta)
    =
    \log\widehat Z_1(\beta)
    +
    \log\widehat Z_2(\beta).
    \tag{108}
\]
The script should print
\[
    \widehatell_2(\beta_0),
    \qquad
    \partial_\beta\widehatell_2(\beta_0),
    \qquad
    D(h),
    \qquad
    |D(h)-\partial_\beta\widehatell_2(\beta_0)|.
    \tag{109}
\]
It should also print the branch manifest or manifest hash used for
\(\beta_0\), \(\beta_0+h\), and \(\beta_0-h\).  The test passes only if those
branch identifiers are identical and the derivative error has a stable
decreasing window before numerical roundoff dominates.

\section{Closing Summary}

The fixed-branch gradient is built from five elementary ideas:
\[
\begin{array}{ll}
\text{normalizer derivative} & \partial\log Z=\dot Z/Z,\\[2mm]
\text{squared density derivative} & \partial(\phi^2)=2\phi\dot\phi,\\[2mm]
\text{mass contraction derivative} & \text{differentiate every factor once},\\[2mm]
\text{linear solve derivative} & N\dot g=\dot d-\dot N g,\\[2mm]
\text{carried filter derivative} & \partial(a/Z)=\dot a/Z-a\dot Z/Z^2.
\end{array}
\tag{110}
\]
The tensor-train notation makes these ideas look large, but it does not change
their nature.  Once the branch is fixed, the forward pass defines one scalar.
The derivative pass differentiates that scalar.  What remains for empirical
work is to test whether the approximation is accurate enough, stable enough,
and useful enough for the target scientific model.

\end{document}
